电力调度中心要实时推断电网各节点的电压相位。节点上百，若把所有节点的状态当成一个整体来描述，参数量随节点数指数增长，枚举根本无从谈起。可电网的物理连接是稀疏的：绝大多数节点之间没有直接耦合，一个节点的状态在给定它的直接邻居之后，与远处节点再无瓜葛。把这句话画成一张图，联合分布就被拆成了一堆只涉及少数变量的局部因子。

这是本单元的全部内容：用条件独立把大分布压成小因子。压缩带来的好处是参数量与计算量一起下降，代价是那些声明必须为真——图上缺的每一条边都是一个断言，断言错了程序不会报错，只会让推断结果悄悄偏掉。

\begin{center}
\begin{tikzpicture}[>=stealth,line width=0.7pt]
\draw (0,0) rectangle (3.1,1.0);
\draw (3.9,0) rectangle (7.0,1.0);
\draw (7.8,0) rectangle (10.9,1.0);
\node[font=\small] at (1.55,0.68) {静态图};
\node[font=\small] at (1.55,0.28) {第一、二篇};
\node[font=\small] at (5.45,0.68) {时间链};
\node[font=\small] at (5.45,0.28) {第三、四篇};
\node[font=\small] at (9.35,0.68) {选型对照};
\node[font=\small] at (9.35,0.28) {第五篇};
\draw[->] (3.1,0.5) -- (3.9,0.5);
\draw[->] (7.0,0.5) -- (7.8,0.5);
\node[font=\small,anchor=north] at (1.55,-0.35) {哪些边可以不画};
\node[font=\small,anchor=north] at (5.45,-0.35) {链结构让推断变线性};
\node[font=\small,anchor=north] at (9.35,-0.35) {谁在条件线右边};
\draw[black!45] (0,-1.35) -- (10.9,-1.35);
\node[font=\small,anchor=north] at (5.45,-1.45) {贯穿三段的同一个量：条件独立};
\end{tikzpicture}
\end{center}

\emph{三段路线共用一条主线，只是把它先后放在静态变量、时间轴和建模选型上。}

第一篇建立共同语言：有向图的分解与 d-分离，无向图的势函数与配分函数，两族图各自表达不了的独立性模式，以及精确推断的代价怎样被树宽而非节点数决定。第二篇把这套语言放进联合 Gauss 的情形，那里协方差的零与精度矩阵的零分别对应边缘独立与条件独立，稀疏结构直接写在矩阵元素上。

第三、四篇沿时间轴展开。Kalman filter 是把 Gauss 条件化沿链反复调用的产物，它顺便交代了可观测性、滤波与平滑的信息窗口差别，以及协方差更新的数值陷阱。条件随机场则把整段输入搬到条件线右边，用全局归一化和前向--后向把 \(K^{T}\) 条路径压成沿边遍历。第五篇把四类序列模型摆在一起，说明它们建模的对象不同，比较强弱之前得先确定要问的是哪类问题。

离散隐状态的入口留在\hyperref[sec:13-Machine-Learning/06-Probabilistic-Models/05]{``HMM 把隐变量排列成时间链''}，本单元不重复讲它的基本设定，只在需要对照时引用。近似推断的两条主路——采样与变分——集中在\hyperref[ch:116]{``采样与近似推断：当目标量写得出却算不出''}；把箭头升级成干预语义则要等到\hyperref[ch:120]{``因果推断：预测世界不等于改变世界''}。

首次阅读建议抓住``分解—独立—消元''这条线：先看图怎样声明局部因子，再区分协方差的零与精度的零，然后观察链结构如何让归一化、边缘概率与解码都落进动态规划。d-分离的完整规则与消息传递的实现细节可以留到第二遍。有一件事从第一遍就要分清——统计图、推断算法与因果图是三样东西，箭头与缺边只在所声明的假设下才有相应含义。

本章篇目按概念依赖排列；若从具体困惑进入，也可以直接打开最接近的一篇，再沿篇内引用回补。

\begin{itemize}
\tightlist
\item \hyperref[sec:13-Machine-Learning/07-Graphical-and-Sequential-Models/01]{``图结构把联合分布拆成局部因子''}；
\item \hyperref[sec:13-Machine-Learning/07-Graphical-and-Sequential-Models/02]{``Gaussian 图模型用精度矩阵编码条件独立''}；
\item \hyperref[sec:13-Machine-Learning/07-Graphical-and-Sequential-Models/03]{``Kalman filter 是动态系统上的 Gaussian 条件化''}；
\item \hyperref[sec:13-Machine-Learning/07-Graphical-and-Sequential-Models/04]{``CRF 对整条标签序列做条件归一化''}；
\item \hyperref[sec:13-Machine-Learning/07-Graphical-and-Sequential-Models/05]{``四类序列模型回答不同问题''}。
\end{itemize}

读完这五篇，你手上会多一种检查清单：拿到一张图模型，先问它声明了哪些条件独立，再问这些声明凭什么成立。
